\section{Tooling}

The Python environment was newer than the specification assumed, running Python
3.14. This was a real risk for the frontend, since a compiler that cannot import
its own dependencies is stuck at the first phase. The risk was retired early by
confirming that numpy, nanobind, onnx, onnxruntime, and a CPU build of torch all
publish wheels for this interpreter, so the frontend proceeded on the existing
environment rather than a downgraded one.

On the C++ side, LLVM builds with exceptions and run time type information
disabled. The first version of the binary decoder used exceptions to signal a
malformed file, which did not compile. Rather than fight the build flags per
target, the decoder was refactored to return an optional, which is a cleaner
interface anyway and bounds any corrupt length field so a bad file cannot request
a huge allocation.

A smaller but recurring friction was invoking the Linux toolchain from the Windows
side. Multi line commands passed through the shell bridge had their quoting
mangled, so the reliable pattern became writing a short script to a scratch
location and running that script, rather than passing complex command lines
inline.
